\chapter{Microscopic information, kinetic limits, and fluid closure}
\label{ch:microscopic-fluid-limits}

\reviewstatus{The 249 timed English transcript segments of the requested
braintruffle video, from 0:00 through 32:36, were reviewed in the public
YouTube transcript interface.  The timestamped source audit records precisely
which accompanying papers, lecture decks, and book sections were accessible.
The full main texts of Stodolna, Vozzi, Itatani, Allen, the abridged English
J\"onsson article, and all six requested Volkov decks were reviewed.  Several
other requested originals were accessible only as abstracts or metadata;
these are not represented as full-text reviews.  The arguments proved below
are ordinary mathematical proofs.  Their presence here does not assert a
Lean formalization of the continuum theories.}

\section{What the microscopic perspective establishes}

The video organizes fluid simulation around successive reductions in the
information retained about a physical system: a quantum state, molecular
positions and velocities, a velocity distribution, and a few macroscopic
fields \citep{microBraintruffle2021}.  This is a productive organizing principle
provided that each reduction is supplied with its own assumptions.  An
exact change of coordinates, a statistical marginal, an asymptotic limit,
and an empirical constitutive law are mathematically different operations.
The distinction determines which conclusions can be carried from one level
to another.

For example, expanding a wave function in a complete eigenbasis retains
all information.  Keeping only finitely many modes discards information.
Taking the partial trace of a many-body density operator discards
correlations.  Replacing a quantum nuclear state by a classical trajectory
requires a localization or semiclassical approximation.  Replacing an
$N$-particle classical distribution by a one-particle distribution requires
a closure of its correlation hierarchy.  Replacing that distribution by
density, velocity, and temperature requires another closure.  None of these
steps follows merely from a count of variables.

The central mathematical lesson is therefore a sequence of conditional
statements.  The Born--Huang expansion identifies the terms discarded by an
adiabatic approximation.  Liouville's equation identifies the correlations
discarded by a kinetic approximation.  Collision invariants identify the
moments retained by a fluid approximation.  An entropy inequality restricts
admissible closures but does not determine all their coefficients.  A
hydrodynamic-limit theorem supplies compactness and convergence information
that a formal moment calculation does not supply.  The structural results
in Chapter~\ref{ch:fluid-structures} and the quantum algorithms in
Chapter~\ref{ch:quantum-fluid} fit into these particular stages.

\subsection{Five distinct mathematical state spaces}

For $N$ distinguishable spinless particles in three dimensions, a pure
quantum state is a unit vector in $L^2(\R^{3N})$, up to an overall phase.
An $N$-particle classical microstate is a point of $\R^{6N}$.  An uncertain
classical state is instead a probability measure on that same phase space;
when it has a density, that density is a function of $6N$ variables.  Thus
the video's contrast between a wave function of $3N$ coordinates and a
classical list of $6N$ numbers compares a general quantum state with a
single classical microstate.  A comparison of statistical descriptions must
also include classical ensemble distributions.

The one-particle kinetic density $f(t,x,v)$ is a function on
$\R^3_x\times\R^3_v$.  In the unit-mass convention used below, its basic
moments are
\begin{equation}
 \rho=\int f\,dv,\qquad
 \rho u=\int vf\,dv,\qquad
 3\rho\theta=\int |v-u|^2f\,dv.
 \label{eq:micro-basic-moments}
\end{equation}
Here $\theta=k_BT/m$ in physical units.  The fluid variables
$(\rho,u,\theta)$ contain five scalar fields.  They generally do not
determine $f$, its pressure tensor, or its heat flux.

In the quantum setting, mixed states are positive trace-one operators
$\varrho$.  A reduced one-particle density matrix and a classical kinetic
density are different kinds of objects.  A Wigner transform can relate
them in suitable semiclassical regimes, but a generic Wigner function is
not a nonnegative probability density.  A classical probabilistic
interpretation therefore cannot be inserted without a positivity or
limiting argument.

\section{Measurement, phase, and the limits of tomography}

The experiments cited by the video support wave interference and
model-dependent reconstruction of electronic structure.  J\"onsson's
multi-slit experiment compares electron diffraction with interference
predictions \citep{microJonsson1974}.  Stodolna and collaborators image
nodal structures in an outgoing electron distribution associated with
hydrogen Stark states; the detector distribution is related to a
propagated wave by the apparatus and its model
\citep{microStodolna2013}.  Itatani's molecular-orbital tomography uses
high-harmonic emission, alignment information, a reference atom, and a
recollision approximation \citep{microItatani2004}.  Vozzi's generalized
tomography improves phase recovery and the treatment of Coulomb
distortion while retaining finite-bandwidth and reconstruction assumptions
\citep{microVozzi2011}.

These are not measurements of a particle trajectory or unrestricted
tomography of an arbitrary many-electron wave function.  In particular, a
single orbital reconstruction does not by itself determine every
many-particle correlation.  The distinction can be expressed without
depending on an interpretation of quantum mechanics.

\begin{proposition}[Finite-dimensional informational completeness]
\label{prop:micro-tomography}
Let $E_1,\ldots,E_s$ be positive operators on $\C^d$ satisfying
$\sum_iE_i=I$.  The probabilities
$p_i=\Tr(E_i\varrho)$ uniquely determine every density operator
$\varrho$ if and only if the real linear span of the $E_i$ is the
space of Hermitian $d\times d$ matrices.
\end{proposition}
\begin{proof}
If the span is the full Hermitian space and two states give the same
probabilities, their difference $X$ is orthogonal to every Hermitian
matrix in the inner product $\Tr(XY)$.  Taking $Y=X$ gives
$\Tr(X^2)=0$, hence $X=0$.

Conversely, if the span is proper, choose a nonzero Hermitian matrix $X$
orthogonal to it.  Since $I=\sum_iE_i$ belongs to the span,
$\Tr X=0$.  For sufficiently small $a>0$, both
$\varrho_\pm=I/d\pm aX$ are positive trace-one matrices.  They are
distinct and satisfy $\Tr(E_i\varrho_+)=\Tr(E_i\varrho_-)$ for
every $i$.
\end{proof}

Position probabilities alone fail this criterion.  The qubit vectors
$(|0\rangle+|1\rangle)/\sqrt2$ and
$(|0\rangle-|1\rangle)/\sqrt2$ have identical computational-basis
probabilities and different relative phase.  Similarly,
$\psi(x)$ and $e^{i\alpha(x)}\psi(x)$ have the same $|\psi(x)|^2$
where both define admissible states, but may have different currents and
future dynamics.  The Born density for a complex wave function is
$|\psi|^2$, not its algebraic square.

\begin{proposition}[Local operations do not transmit a signal through a marginal]
\label{prop:micro-nosignalling}
Let $\varrho_{AB}$ be a density operator on a finite tensor product.
For operators $K_j$ on $B$ satisfying $\sum_jK_j^*K_j=I_B$, define
\[
 \widetilde\varrho_{AB}
 =\sum_j(I_A\otimes K_j)\varrho_{AB}(I_A\otimes K_j^*).
\]
Then $\Tr_B\widetilde\varrho_{AB}=\Tr_B\varrho_{AB}$.
\end{proposition}
\begin{proof}
For every observable $A$ on the first factor, cyclicity of the trace gives
\[
 \Tr((A\otimes I)\widetilde\varrho)
 =\sum_j\Tr((A\otimes K_j^*K_j)\varrho)
 =\Tr((A\otimes I)\varrho).
\]
Equality against every $A$ determines the reduced operator.
\end{proof}

Conditioning on a recorded outcome can change a conditional state.  The
outcome record must then be communicated before the other observer can
condition on it.  Bell correlations consequently do not justify an
assertion of controllable instantaneous communication.  This qualification
is needed when using the introductory discussion in van Dommelen's
book \citep{microVanDommelen}.

\subsection{What the multi-slit formula proves}

In the Fraunhofer approximation, suppose $s$ equally illuminated identical
slits contribute a common amplitude $a(\vartheta)$ and a phase increment
$\delta=2\pi d\sin\vartheta/\lambda$ between neighboring slits.
Linearity gives
\[
 A(\vartheta)=a(\vartheta)\sum_{j=0}^{s-1}e^{ij\delta}.
\]
The geometric-series identity and continuity at removable singularities
give
\begin{equation}
 |A(\vartheta)|^2
 =|a(\vartheta)|^2
 \frac{\sin^2(s\delta/2)}{\sin^2(\delta/2)}.
 \label{eq:micro-slits}
\end{equation}
For a uniformly illuminated slit of width $b$, the Fourier integral
$\int_{-b/2}^{b/2}e^{2\pi i x\sin\vartheta/\lambda}\,dx$
supplies a squared sinc envelope.  Equation~\eqref{eq:micro-slits}
therefore proves the interference factor under stated coherence,
illumination, and far-field assumptions.  It does not remove the need to
model the source, finite detector resolution, or partial coherence.
The reviewed 1974 article explicitly identifies itself as an abridged
translation; it cannot stand in for a complete review of the longer
1961 German original \citep{microJonsson1961,microJonsson1974}.

\section{Exterior powers are a physical state constraint}

The project's exterior algebra is directly relevant to identical
fermions.  Let $V$ be a one-particle space including spin.  The
antisymmetric $N$-particle subspace is $\bigwedge^N V$, identified
with the antisymmetric tensors inside $V^{\otimes N}$.  Exchanging two
complete spin-space labels changes a fermionic state by a sign.  Equality
of position probabilities under an exchange, by itself, does not specify
that sign or the spin dependence.

\begin{proposition}[Slater normalization and dimension]
\label{prop:micro-slater}
If $e_1,\ldots,e_M$ are an orthonormal basis of the finite-dimensional
one-particle space $V$, then for
$1\le i_1<\cdots<i_N\le M$ the tensors
\[
 \Psi_{i_1\cdots i_N}
 =\frac1{\sqrt{N!}}\sum_{\sigma\in S_N}
   \operatorname{sgn}(\sigma)
   e_{i_{\sigma(1)}}\otimes\cdots\otimes e_{i_{\sigma(N)}}
\]
form an orthonormal basis of $\bigwedge^N V$.  In particular,
$\dim\bigwedge^N\C^M=\binom MN$, with the dimension zero for
$N>M$.
\end{proposition}
\begin{proof}
For fixed distinct indices the $N!$ elementary tensors in the sum are
mutually orthonormal, so the squared norm is $N!/N!=1$.  Different
index sets have disjoint elementary-tensor supports.  Antisymmetry forces
the coefficient of any tensor with a repeated index to vanish, and relates
all coefficients belonging to permutations of a distinct index set by
the permutation sign.  Thus every antisymmetric tensor is a linear
combination of the displayed vectors.
\end{proof}

In position representation this is the normalized Slater determinant.
The sign representation of the symmetric group selects the Young diagram
$(1^N)$, while symmetric bosonic tensors select $(N)$.  More general
Schur functors organize symmetry types of tensor features and multi-index
responses.  This is a precise use of the algebraic material developed
earlier in the monograph: a representation specifies allowed states or
tensor components.  It does not, on its own, furnish a bound for a nonlinear
fluid equation.

\subsection{Number states and coherent states}

Mathur's accessible introductory modules are useful for interference,
tensor products, and the distinction between classical amplitude and
measurement probabilities \citep{microMathur}.  One teaching
identification needs correction: a state with a fixed number of identical
bosons in one mode is a number state, and is not a coherent state.
Let $a|n\rangle=\sqrt n\,|n-1\rangle$.  A coherent state is
\begin{equation}
 |\alpha\rangle=e^{-|\alpha|^2/2}
       \sum_{n=0}^{\infty}\frac{\alpha^n}{\sqrt{n!}}|n\rangle.
 \label{eq:micro-coherent}
\end{equation}
The exponential series shows that its squared norm is one.  Shifting
the summation index shows $a|\alpha\rangle=\alpha|\alpha\rangle$.
Its number probabilities are
$e^{-|\alpha|^2}|\alpha|^{2n}/n!$, whose mean and variance are both
$|\alpha|^2$.  A fixed $|n\rangle$ has number variance zero and
$\langle a\rangle=0$.  Quantization of the number observable therefore
does not discretize every complex amplitude or every field expectation.
This matters when deciding what a quantum neural model's input amplitudes
and measured outputs represent.

\section{Born--Huang expansion before Born--Oppenheimer truncation}

The video's move from electronic wave functions to molecular potential
energy surfaces is best made precise by first writing the coupling terms
that are later neglected.  The treatment in van Dommelen's Section 9.2
and derivation D.51 distinguishes electronic and nuclear quantum dynamics;
Allen's molecular-dynamics notes then describe classical motion on a chosen
potential \citep{microVanDommelen,microAllen2004}.

For an exact, domain-transparent derivation, take a finite electronic
space $\C^r$, nuclear coordinate $R\in\R^d$, and a $C^2$ family of
Hermitian matrices $H_e(R)$.  Suppose an orthonormal eigenbasis
$\phi_n(R)$ can be chosen $C^2$ on the open region being considered.
For $C^2$ nuclear coefficient functions $\chi_n$, set
\[
 H=-\frac{\hbar^2}{2M}\Delta_R+H_e(R),\qquad
 \Psi(R)=\sum_{n=1}^r\chi_n(R)\phi_n(R).
\]
Scalar nuclear repulsion can be included in $H_e$.  The corresponding
continuum-electronic formulas require common operator domains and the
regularity needed to justify these derivatives; the finite statement
does not assert those analytic hypotheses automatically.

\begin{proposition}[Exact coupled nuclear equations]
\label{prop:micro-bornhuang}
Set
\[
 \begin{aligned}
 d_{mn}&=\langle\phi_m,\nabla_R\phi_n\rangle,\\
 \tau_{mn}&=\langle\phi_m,\Delta_R\phi_n\rangle.
 \end{aligned}
\]
Then $H\Psi=E\Psi$ is equivalent to
\begin{equation}
 -\frac{\hbar^2}{2M}
  \left(\Delta\chi_m+
        2\sum_n d_{mn}\cdot\nabla\chi_n+
        \sum_n\tau_{mn}\chi_n\right)
  +E_m(R)\chi_m=E\chi_m
 \label{eq:micro-bornhuang}
\end{equation}
for each $m$.  Moreover, $d_{mn}=-\overline{d_{nm}}$ componentwise.
\end{proposition}
\begin{proof}
Apply the product rule to each
$\Delta(\chi_n\phi_n)$ and take its electronic inner product with
$\phi_m$.  Orthonormality gives the first term and diagonalizes
$H_e$.  Differentiating
$\langle\phi_m,\phi_n\rangle=\delta_{mn}$ proves the final identity.
\end{proof}

The usual single-surface approximation retains one electronic sector and
controls or omits its couplings to others.  Which diagonal geometric
terms are retained must be specified.  Rephasing
$\phi_n(R)\mapsto e^{i\gamma_n(R)}\phi_n(R)$ also rephases
$\chi_n$ and changes diagonal derivative couplings; a claim that they
vanish everywhere needs a gauge and topology argument.

\begin{proposition}[Energy gradients and the gap denominator]
\label{prop:micro-hellmann}
Let $E_n(R)$ be a simple normalized eigenpair of the finite matrix family.
For a coordinate derivative $\partial_j$,
\[
 \partial_j E_n
 =\langle\phi_n,(\partial_jH_e)\phi_n\rangle,
 \qquad
 \langle\phi_m,\partial_j\phi_n\rangle
 =\frac{\langle\phi_m,(\partial_jH_e)\phi_n\rangle}
        {E_n-E_m}\quad(m\ne n).
\]
If $|E_n-E_m|\ge\gamma>0$ for all $m\ne n$, then
\[
 \|(I-|\phi_n\rangle\langle\phi_n|)\partial_j\phi_n\|
 \le \gamma^{-1}\|\partial_jH_e\|.
\]
\end{proposition}
\begin{proof}
Differentiate $H_e\phi_n=E_n\phi_n$ and take inner products
with $\phi_n$ and with $\phi_m$.  The terms containing
$\partial_j\phi_n$ cancel in the first case.  In the second they
give the difference $E_n-E_m$.  Squaring the off-diagonal coefficients
and summing over the orthonormal eigenbasis gives the stated norm bound.
\end{proof}

Large nuclear mass reduces the prefactor in
\eqref{eq:micro-bornhuang}; it does not prevent a small electronic gap
from amplifying the derivative coupling.  Near crossings and conical
intersections, a single-surface argument requires additional work.
Likewise, an excited electronic surface may have an adiabatic nuclear
description.  The Born--Oppenheimer approximation is not a theorem that
electrons automatically dissipate into their ground state.

\begin{proposition}[Unitary evolution does not supply ground-state relaxation]
\label{prop:micro-no-relaxation}
For a time-independent finite Hermitian matrix $H$ and a spectral
projector $P$, the population
\[
 \Tr(Pe^{-itH/\hbar}\varrho_0e^{itH/\hbar})
\]
is constant.
\end{proposition}
\begin{proof}
The projector commutes with both exponentials.  Cyclicity of the trace
reduces the expression to $\Tr(P\varrho_0)$.
\end{proof}

A bath, emitted radiation, collisions with other degrees of freedom, or
a specified open-system approximation can generate relaxation.
Adiabatic following and relaxation are distinct physical mechanisms.
The video at 21:37--21:53 compresses them into a single explanation;
equations~\eqref{eq:micro-bornhuang} and
Proposition~\ref{prop:micro-no-relaxation} state what must be separated.

\section{When a quantum packet behaves approximately as a particle}

For
$H=P^2/(2M)+V(X)$, the formal Ehrenfest identities are
\begin{equation}
 \frac{d}{dt}\langle X\rangle=\frac{\langle P\rangle}{M},
 \qquad
 \frac{d}{dt}\langle P\rangle=-\langle\nabla V(X)\rangle.
 \label{eq:micro-ehrenfest}
\end{equation}
They follow from $d\langle A\rangle/dt=(i/\hbar)\langle[H,A]\rangle$
on a common invariant domain.  For instance, they hold for sufficiently
regular rapidly decaying solutions for which all integrations by parts
and differentiated moments are justified.  They are not yet Newton's
equations for the mean, because
$\langle\nabla V(X)\rangle$ need not equal
$\nabla V(\langle X\rangle)$.

\begin{proposition}[Localization controls the mean-force error]
\label{prop:micro-ehrenfest-error}
If $\nabla V$ is $L$-Lipschitz and the position distribution has a
finite second moment, then
\[
 \left|\langle\nabla V(X)\rangle-
       \nabla V(\langle X\rangle)\right|
 \le L\left\langle|X-\langle X\rangle|^2\right\rangle^{1/2}.
\]
\end{proposition}
\begin{proof}
Apply the Lipschitz inequality inside the position integral, then the
triangle inequality and Cauchy--Schwarz for its probability measure.
\end{proof}

This gives a quantitative condition rather than a universal particle-size
threshold.  To turn it into a trajectory approximation on a time interval,
one must control the packet width throughout that interval and estimate
the resulting perturbation of the classical flow.

For a free particle in one dimension with $\sigma>0$ and
\[
 \psi_0(x)=(2\pi\sigma^2)^{-1/4}
             e^{-x^2/(4\sigma^2)},
\]
the exact Heisenberg identity $X(t)=X+tP/M$ gives
\begin{equation}
 \operatorname{Var}(X(t))
 =\sigma^2+\frac{\hbar^2t^2}{4M^2\sigma^2}.
 \label{eq:micro-gaussian}
\end{equation}
Indeed $\operatorname{Var}X=\sigma^2$,
$\operatorname{Var}P=\hbar^2/(4\sigma^2)$, and the initial
symmetrized covariance vanishes.  Increasing $M$ slows free spreading
for a fixed packet, but \eqref{eq:micro-gaussian} still grows with time.
Heavy particles can also exhibit interference.  A suitable scale
comparison involves action, wavelength, temperature, environmental
coupling, and the observable being approximated.

\subsection{One-dimensional Coulomb models require a domain}

The video explicitly uses a softened one-dimensional potential to explain
electronic forces.  Such a model is pedagogically useful, but it is a
different operator from the singular expression
\[
 -\frac{\hbar^2}{2m}\frac{d^2}{dx^2}-\frac{\kappa}{|x|}.
\]
A differential expression does not specify a self-adjoint Hamiltonian
until its domain and behavior at the singular point are given.
The reviewed operator-domain analysis by de Oliveira and Verri finds a
family of self-adjoint realizations for the one-dimensional operator
initially defined away from the origin
\citep{microOliveiraVerri2009}.  We use that domain distinction, not
an unqualified assertion about the spectrum of every regularization.
The requested Loudon papers are separately registered as unavailable
full texts in this review \citep{microLoudon1959,microLoudon2016}.

The limit of a family of softened potentials must specify the topology
of operator convergence and any renormalization of energy.  It cannot
be inferred merely by drawing potentials that approach $-\kappa/|x|$
away from zero.  This is the same structural issue that recurs in
continuum limits of graph Laplacians: a formula for coefficients must be
accompanied by domains, boundary conditions, and convergence estimates.

\section{Molecular dynamics and conserved microscopic information}

A classical molecular Hamiltonian has the form
\begin{equation}
 H(q,p)=\sum_{i=1}^N\frac{|p_i|^2}{2m_i}+U(q),\qquad
 \dot q_i=\frac{p_i}{m_i},\quad
 \dot p_i=-\nabla_{q_i}U.
 \label{eq:micro-md}
\end{equation}
The force is the negative gradient of potential energy.  General
molecular force fields include angular, torsional, many-body, constrained,
or reactive terms; they need not be sums of pairwise distance functions.
Allen's notes explain these modeling choices, symplectic integrators,
constraints, thermostats, and statistical errors
\citep{microAllen2004}.  Their scope is much wider than the two-particle
illustration in the video.

\begin{proposition}[Liouville volume and fine-grained entropy]
\label{prop:micro-liouville}
Let $H\in C^2(\R^{2n})$ have a smooth Hamiltonian flow $\Phi_t$
on the interval considered.  Then $\det D\Phi_t=1$.
If a smooth probability density evolves by
$f_t=f_0\circ\Phi_{-t}$ and the entropy integral is absolutely
integrable, then
\[
 \int f_t\log f_t\,dq\,dp
 =\int f_0\log f_0\,dq\,dp.
\]
\end{proposition}
\begin{proof}
For the Hamiltonian vector field
$X_H=(\partial_pH,-\partial_qH)$, equality of mixed derivatives
gives
$\operatorname{div}X_H=\sum_i(\partial_{q_i}\partial_{p_i}H-
\partial_{p_i}\partial_{q_i}H)=0$.
Jacobi's determinant identity for the flow derivative yields
$\frac d{dt}\det D\Phi_t=(\operatorname{div}X_H)\circ\Phi_t
\det D\Phi_t=0$.  Its initial value is one.  A change of variables
in the entropy integral proves the second assertion.
\end{proof}

Thus a Hamiltonian description preserves fine-grained information.
Entropy growth in a reduced description is not the destruction of
information by an invertible microscopic flow.  It involves a
marginalization, an environment, a limiting operation, or a restriction
on what is resolved.  Quantum unitary conjugation similarly preserves
the spectrum of a density operator and hence its von Neumann entropy
when that entropy is defined.  A reduced subsystem can have changing
entropy because its correlations with the omitted subsystem change.

\subsection{A useful exact statement about velocity Verlet}

For separable $H=T(p)+U(q)$, the kick
$(q,p)\mapsto(q,p-h\nabla U(q))$ and drift
$(q,p)\mapsto(q+h\nabla T(p),p)$ are exact Hamiltonian subflows.
Their symmetric composition, kick by $h/2$, drift by $h$, kick by
$h/2$, is symplectic.  This proves preservation of the symplectic form;
it does not prove exact preservation of the original $H$.

For the scalar oscillator $H=(p^2+\omega^2q^2)/2$, the update is
\[
 q'=\left(1-\frac{h^2\omega^2}{2}\right)q+hp,
 \quad
 p'=\left(1-\frac{h^2\omega^2}{2}\right)p
       -h\omega^2\left(1-\frac{h^2\omega^2}{4}\right)q.
\]
Direct expansion proves that it preserves
\begin{equation}
 \widetilde H_h(q,p)
 =\frac12p^2+\frac12\omega^2
       \left(1-\frac{h^2\omega^2}{4}\right)q^2.
 \label{eq:micro-verlet-invariant}
\end{equation}
This quadratic form is positive definite when $0<h\omega<2$.
It is a valid special-case modified invariant.  In a general nonlinear
system, a formal or asymptotic modified Hamiltonian from backward error
analysis must not be promoted to an exact global invariant without the
additional theorem needed for that claim.

\section{Marginals, correlations, and the kinetic approximation}

For a smooth symmetric $N$-particle density $F_N$, let $F_N^{(s)}$
denote its $s$-particle marginal.  With unit masses and a pair force
$K(x)=-\nabla\Phi(x)$, integrating the Liouville equation over all
but particle 1 gives
\begin{equation}
 \partial_tF_N^{(1)}+v\cdot\nabla_xF_N^{(1)}
 +(N-1)\nabla_v\cdot
   \int K(x-y)F_N^{(2)}(x,v,y,w)\,dy\,dw=0,
 \label{eq:micro-bbgky-first}
\end{equation}
assuming the boundary terms vanish.  The two-particle marginal appears
because the force on a particle depends on others.  Its own equation
involves the three-particle marginal.  This is the BBGKY hierarchy.

\begin{example}[A first marginal cannot determine a pair correlation]
\label{ex:micro-marginal}
On $\{0,1\}^2$, let $P$ give mass $1/2$ each to $(0,0)$ and
$(1,1)$, and let $Q$ give mass $1/2$ each to $(0,1)$ and $(1,0)$.
Both first and second one-coordinate marginals are uniform.  Nevertheless,
$P(X=Y)=1$ and $Q(X=Y)=0$.  Thus no function of the one-coordinate
marginals alone can reconstruct all pair observables.
\end{example}

Replacing $F_N^{(2)}$ by a product is a substantive assumption.
For dilute collisions it is the incoming, appropriately scaled
correlations that matter.  Hard-sphere exclusion already prevents exact
independence at fixed particle diameter.  The word ``chaos'' in
propagation of chaos denotes an asymptotic factorization property of
fixed-order marginals; it is not interchangeable with deterministic
sensitivity to initial conditions.

\subsection{A rigorous classical-to-kinetic bridge and its time interval}

In the hard-sphere Boltzmann--Grad scaling in dimension $d\ge2$,
\[
 N\varepsilon^{d-1}=1,
 \qquad N\varepsilon^d=\varepsilon\longrightarrow0.
\]
The first relation retains an order-one collision scale, while the
second sends the occupied volume fraction to zero.  Gallagher,
Saint-Raymond, and Texier give a detailed derivation for hard spheres
and suitable short-range repulsive potentials
\citep{microGallagher2013}.  Their hard-sphere theorem propagates
admissible asymptotic initial factorization on a positive short time
interval, with convergence of velocity-averaged observables locally
away from positional diagonals.  It uses weighted uniform bounds on
the hierarchy and quantitative control of recollisions.

The reviewed source's introductory statement assumes a continuous
probability density $f_0$ with a Gaussian velocity bound
\[
 \sup_{x,v}f_0(x,v)e^{\beta|v|^2/2}<\infty
 \quad\text{for some }\beta>0.
\]
Its precise Theorem 8 includes the admissibility and hierarchy-space
hypotheses.  The proof does not assert convergence for arbitrary
finite-$N$ initial data, arbitrarily long times, every potential, or
every boundary geometry.  This chapter reviewed the introduction,
main convergence statement, and concluding scope discussion; it does
not claim to reproduce or independently verify the monograph's complete
recollision proof.

This time restriction has a practical logical consequence.  A short-time
particle-to-Boltzmann theorem and a long macroscopic-time
Boltzmann-to-fluid theorem cannot simply be concatenated.  The former
must hold uniformly on the time scales required by the latter, or a
separate joint-limit theorem must be supplied.  No amount of formal
diagram composition replaces that analytic requirement.

\section{The collision algebra and its five invariants}

We now prove the algebraic facts behind the simplest elastic collision
operator.  They are useful both for conservative simulation and for
checking learned collision models.  Set $g=v-v_*$, let
$n\in S^{d-1}$, and define
\begin{equation}
 v'=v-(g\cdot n)n,\qquad
 v_*'=v_*+(g\cdot n)n.
 \label{eq:micro-collision}
\end{equation}

\begin{proposition}[Elastic reflection]
\label{prop:micro-collision}
The map \eqref{eq:micro-collision} is an involution with absolute
Jacobian determinant one.  It preserves total momentum and kinetic
energy:
\[
 v'+v_*'=v+v_*,\qquad
 |v'|^2+|v_*'|^2=|v|^2+|v_*|^2.
\]
\end{proposition}
\begin{proof}
The center velocity $V=(v+v_*)/2$ is unchanged and
$g'=g-2(g\cdot n)n$.  This is an orthogonal reflection, so
$|g'|=|g|$, its square is the identity, and its determinant has
absolute value one.  Since
$|v|^2+|v_*|^2=2|V|^2+|g|^2/2$, all assertions follow after the
invertible linear change from $(v,v_*)$ to $(V,g)$.
\end{proof}

A scalar function $\phi$ is a collision invariant when
$\phi(v)+\phi(v_*)=\phi(v')+\phi(v_*')$ for every elastic
collision.  The functions $1,v_1,\ldots,v_d,|v|^2$ are immediately
invariant.  The following proof makes the converse precise under a
simple regularity assumption.

\begin{proposition}[Classification of smooth collision invariants]
\label{prop:micro-invariants}
For $d\ge2$, every $C^2$ collision invariant has the form
\[
 \phi(v)=a+b\cdot v+c|v|^2.
\]
Consequently the scalar invariant space has dimension $d+2$, hence
five in three dimensions.
\end{proposition}
\begin{proof}
For fixed $V$ and radius $r$, collisions can send $g$ to any vector
of the same length: reflection in a suitable hyperplane does so for
two given vectors.  Hence
$\phi(V+r\omega)+\phi(V-r\omega)$ is independent of the unit
vector $\omega$.  Its second-order expansion at $r=0$ implies
that $\omega^TD^2\phi(V)\omega$ is constant on the unit sphere.
Taking coordinate vectors and their pairwise sums shows
$D^2\phi(V)=\lambda(V)I$.

The mixed second derivatives are zero, so $\partial_i\phi(v)$
depends only on $v_i$.  The equality of all diagonal second
derivatives then says that a function of $v_i$ alone equals a function
of $v_j$ alone for independently variable $v_i,v_j$.  Since $d\ge2$,
each is the same constant.  Integrating twice yields the displayed
quadratic polynomial.  The converse follows from
Proposition~\ref{prop:micro-collision}.
\end{proof}

Some introductory notes group the invariants into three physical
categories: particle number, momentum, and energy.  Momentum is a vector,
so those categories contain five independent scalar invariants in
three dimensions, not three scalar equations.

\subsection{A collision kernel contains more information than a diameter}

For a central classical interaction, the differential cross section is
obtained from the map between impact parameter $b$ and scattering
angle $\chi$.  On a branch where the map is invertible,
\[
 \frac{d\sigma}{d\Omega}
 =\frac{b}{\sin\chi}\left|\frac{db}{d\chi}\right|.
\]
If multiple branches produce the same angle, their contributions must
be added.  Angular singularities and grazing collisions require further
analysis.  The hard-sphere model and the variable hard-sphere model
used in Volkov's notes do not reproduce every angular distribution of
a general potential \citep{microVolkov02}.  Matching a viscosity
law or a mixing statistic is a calibration of particular observables,
not an equality of the underlying scattering laws.

Internal excitation and chemical reactions add variables, transitions,
and conservation relations.  Parker's requested paper concerns
rotational and vibrational relaxation, while Valentini and collaborators'
requested study concerns nitrogen dissociation on an ab initio potential
surface \citep{microParker1959,microValentini2015}.  Only their
accessible abstracts and metadata were reviewed here.  Their reported
models must not be treated as fully audited derivations or as evidence
that every gas can be represented by an elastic monatomic kernel.

\section{The weak collision identity and entropy production}

Let $B(g,n)\ge0$ be a collision kernel invariant under exchange of the
particles and under the pre/post-collision reflection.  Define
\begin{equation}
 Q(f,f)(v)=\int_{\R^d}\int_{S^{d-1}}
       B(v-v_*,n)(f'f_*'-ff_*)\,dn\,dv_*.
 \label{eq:micro-boltzmann-Q}
\end{equation}
Here all four values of $f$ have the same $(t,x)$, when those
variables are present.  The next statements assume absolute
integrability of the displayed integrals, so that Fubini's theorem
and the changes of variables are legitimate.

\begin{proposition}[Weak collision identity]
\label{prop:micro-weak-Q}
For an admissible test function $\phi$,
\begin{align}
 \int\phi Q(f,f)\,dv
 &=\frac14\iiint B(f'f_*'-ff_*)
       (\phi+\phi_* -\phi'-\phi_*')\,dn\,dv_*\,dv
       \label{eq:micro-weak-Q}\\
 &=\frac12\iiint Bff_*
       (\phi'+\phi_*'-\phi-\phi_*)\,dn\,dv_*\,dv.
       \nonumber
\end{align}
\end{proposition}
\begin{proof}
Exchange $v$ and $v_*$ in the original integral and average the
two expressions.  In a second copy, change variables from
$(v,v_*)$ to $(v',v_*')$.  The measure and $B$ are unchanged,
the gain-minus-loss factor changes sign, and the test functions
acquire primes.  Averaging once more gives the first identity.
Applying the same pre/post change to the gain term gives the second.
\end{proof}

Testing with collision invariants proves the mass, momentum, and energy
collision balances.  Testing with the logarithm supplies a stronger
constraint, subject to the additional integrability needed for that
unbounded test.

\begin{theorem}[Classical homogeneous $H$ theorem under explicit regularity]
\label{thm:micro-H}
Let $f(t,v)>0$ be a differentiable solution of
$\partial_tf=Q(f,f)$ for which $\int f\log f$ is finite,
differentiation under the integral is valid, and all terms in
\eqref{eq:micro-weak-Q} with $\phi=\log f$ are integrable.
Then
\begin{equation}
 \frac{d}{dt}\int f\log f\,dv=-D(f)\le0,
 \quad
 D(f)=\frac14\iiint B(f'f_*'-ff_*)
           \log\frac{f'f_*'}{ff_*}\,dn\,dv_*\,dv.
 \label{eq:micro-H}
\end{equation}
\end{theorem}
\begin{proof}
The derivative equals $\int(1+\log f)Q(f,f)\,dv$.
The constant part vanishes by mass conservation.  Use
\eqref{eq:micro-weak-Q} for the remaining part.  Since
$(a-b)(\log a-\log b)\ge0$ for $a,b>0$, the displayed
entropy production is nonnegative.
\end{proof}

For spatially varying smooth solutions, transport adds the entropy flux
$\nabla_x\cdot\int vf\log f\,dv$.  Boundary entropy fluxes must
then be included.  For weak or renormalized solutions, an entropy
inequality is established by approximation and lower semicontinuity;
one cannot simply insert an unjustified logarithmic test into a weak
equation.

\begin{proposition}[Finite energy does not imply finite positive entropy]
\label{prop:micro-entropy-counterexample}
There is a probability density supported in $[0,1]^3$ with finite
second moment and $\int f\log f=+\infty$.
\end{proposition}
\begin{proof}
Let $c=6/\pi^2$ and choose disjoint measurable sets $A_n$ in
$[0,1]^3$ with volumes
$|A_n|=c e^{-n^2}/n^2$.  Such sets exist because the sum of
their volumes is less than one.  Define $f=e^{n^2}$ on $A_n$
and $f=0$ elsewhere, with $0\log0=0$.  Then
\[
 \int f=c\sum_{n\ge1}n^{-2}=1,
 \qquad \int|v|^2f\,dv\le3,
 \qquad \int f\log f=c\sum_{n\ge1}1=+\infty.
\]
\end{proof}

This corrects a gap in the elementary entropy discussion in the
reviewed Volkov slides \citep{microVolkov03}.  Similarly, integrability
of a continuous density alone does not force pointwise decay to zero
at infinity: separated continuous spikes of height one and summable
widths provide a counterexample.  Decay assumptions used in
integration by parts must therefore be stated, not inferred solely
from normalization.

\subsection{Entropy minimization is not a convergence proof}

\begin{proposition}[Maxwellian minimizes $H$ at fixed first two moments]
\label{prop:micro-maxwell}
Let $f$ be a probability density with mean $u$ and
$\int|v-u|^2f\,dv=d\theta$, where $\theta>0$.  Set
\[
 M_{u,\theta}(v)=(2\pi\theta)^{-d/2}
       e^{-|v-u|^2/(2\theta)}.
\]
Whenever the entropy expressions are defined,
\[
 \int f\log f-\int M_{u,\theta}\log M_{u,\theta}
 =\int f\log\frac{f}{M_{u,\theta}}\ge0.
\]
Equality holds exactly when $f=M_{u,\theta}$ almost everywhere.
\end{proposition}
\begin{proof}
The function $\log M_{u,\theta}$ is a quadratic polynomial in
$v$, so its integral against $f$ agrees with its integral against
$M_{u,\theta}$.  This proves the identity.  With
$r=f/M_{u,\theta}$, use $r\log r-r+1\ge0$ and integrate
against $M_{u,\theta}\,dv$.  Equality in this strictly convex
inequality requires $r=1$ almost everywhere.
\end{proof}

Under a positive kernel, sufficient continuity, and $\log f\in C^2$,
vanishing $D(f)$ forces detailed balance
$f'f_*'=ff_*$.  Proposition~\ref{prop:micro-invariants} applied to
$\log f$ then gives the Maxwellian form, with a negative quadratic
coefficient required by integrability.

These facts characterize entropy minimizers and stationary collision
equilibria.  They do not prove that every trajectory converges to one.
If a scalar function $H(t)$ is nonincreasing and bounded below, it
has a limit, but its derivative need not tend to zero.  For example,
take a nonnegative smooth function $b$ supported in $(-1/4,1/4)$
with $b(0)=1$ and define
\[
 h(t)=\sum_{n\ge1}b(2^n(t-n)),\qquad
 H(t)=\int_t^\infty h(s)\,ds.
\]
The bumps have disjoint supports and summable areas.  Hence $H$ is
smooth, nonnegative, and decreasing to zero, while $H'(n)=-1$.
Even $H'(t)\to0$ alone would not supply compactness of the state
or identify a unique limit.  Convergence-to-equilibrium theorems require
the appropriate coercivity, compactness, integrability, and control of
invariants in the stated regime.

\section{Exact moment balances and the closure problem}

For a smooth unit-mass solution of
\[
 \partial_tf+v\cdot\nabla_xf=Q(f,f),
\]
assume sufficient velocity integrability to justify the following
moments.  Let $c=v-u$ and define
\[
 P=\int c\otimes c\,f\,dv,\quad
 e=\frac1{2\rho}\int|c|^2f\,dv,\quad
 q=\frac12\int|c|^2c f\,dv,
 \quad E=\frac12\rho|u|^2+\rho e.
\]
Here $P$ is a positive semidefinite pressure tensor.  The Cauchy stress
uses the opposite sign for the kinetic contribution; keeping this
convention explicit avoids sign errors in viscosity formulas.

\begin{proposition}[Moment equations before closure]
\label{prop:micro-moments}
Under these hypotheses,
\begin{align}
 \partial_t\rho+\nabla\cdot(\rho u)&=0,
 \label{eq:micro-mass}\\
 \partial_t(\rho u)+\nabla\cdot(\rho u\otimes u+P)&=0,
 \label{eq:micro-momentum}\\
 \partial_tE+\nabla\cdot(Eu+Pu+q)&=0.
 \label{eq:micro-energy}
\end{align}
\end{proposition}
\begin{proof}
Test the kinetic equation with $1,v,|v|^2/2$ and use the collision
invariants.  Since $\int cf\,dv=0$, expansion of
$v=u+c$ gives
$\int v\otimes vf\,dv=\rho u\otimes u+P$.
Likewise,
\[
 \frac12\int |v|^2vf\,dv
 =\left(\frac12\rho|u|^2+\rho e\right)u+Pu+q.
\]
These expansions give all three equations.
\end{proof}

\begin{example}[Equal five moments, different stress]
\label{ex:micro-unclosed}
Let $f_A$ and $f_B$ be centered normalized Gaussian densities on
$\R^3$ with covariance matrices
\[
 A=I,\qquad B=\operatorname{diag}(1+\delta,1-\delta,1),
 \quad0<\delta<1.
\]
They have the same mass, mean, and scalar temperature because
$\tr A=\tr B=3$.  Their pressure tensors are $A$ and $B$,
respectively.  Thus the five fluid moments do not determine the
momentum flux.
\end{example}

A local Maxwellian imposes $P=\rho\theta I$ and $q=0$, giving
the Euler fluxes.  An arbitrary local Maxwellian with space- and
time-dependent parameters is generally not a solution of the full
Boltzmann equation: its collision term vanishes, but its transport
term need not.  Euler equations arise as moment equations under
that closure or as a justified limit, rather than as a proof that
every Euler solution lifts to an exact local-Maxwellian kinetic
solution.  This qualification is stated explicitly in the rigorous
hydrodynamic-limit literature \citep{microGolseSaintRaymond2009}.

\subsection{First-order constitutive laws and scaling}

For a monatomic gas near local equilibrium, a formal
Chapman--Enskog expansion produces the constitutive form
\begin{equation}
 P=pI-2\mu S_0-\zeta(\nabla\cdot u)I,
 \qquad q=-\kappa\nabla T,
 \quad
 S_0=\tfrac12(\nabla u+\nabla u^T)
          -\tfrac13(\nabla\cdot u)I.
 \label{eq:micro-constitutive}
\end{equation}
The bulk coefficient is zero in the simplest monatomic elastic
model.  Internal modes and other models change this conclusion.
The values of the transport coefficients require solving equations
for the linearized collision operator; conservation and symmetry
alone do not fix them.

The Knudsen number is the ratio of a microscopic mean free path to
the macroscopic length scale.  Small Knudsen number, small Mach number,
and fixed Reynolds number specify distinct relationships among
scales.  A formal expansion in one does not silently impose the
others.  A global small value also does not rule out local breakdown
near a shock, boundary layer, or strongly varying rarefied region.
Volkov's first and third decks make the physical scale separation
and the failure of elementary closure in rarefied regimes explicit
\citep{microVolkov01,microVolkov03}.

\subsection{A precise historical hydrodynamic-limit theorem}

The result of Golse and Saint-Raymond that is useful here concerns
the scaled equation on $\R^3_x\times\R^3_v$
\begin{equation}
 \varepsilon^2\partial_tF_\varepsilon
 +\varepsilon v\cdot\nabla_xF_\varepsilon
 =Q(F_\varepsilon,F_\varepsilon),
 \qquad M(v)=(2\pi)^{-3/2}e^{-|v|^2/2}.
 \label{eq:micro-hydro-scale}
\end{equation}
Their collision kernels obey the exchange and reflection symmetries
and hard-cutoff bounds of the form
\[
 0<B(z,n)\le C(1+|z|)^\beta|\widehat z\cdot n|,
 \qquad
 \int_{S^2}B(z,n)\,dn\ge C^{-1}\frac{|z|}{1+|z|},
 \quad0\le\beta\le1,
\]
with the inequalities understood almost everywhere.  Initial data
satisfy
\begin{equation}
 \sup_{\varepsilon>0}\varepsilon^{-2}
 H(F^{\mathrm{in}}_\varepsilon\mid M)<\infty,
 \quad
 H(F\mid M)=\iint\left(F\log\frac FM-F+M\right)\,dx\,dv.
 \label{eq:micro-relative-scaling}
\end{equation}
Their Theorem 2.4 describes limit points of the fluctuations
$g_\varepsilon=(F_\varepsilon/M-1)/\varepsilon$ for the
corresponding renormalized solutions.  In the stated weak local
topology the limit has the form
\[
 g(t,x,v)=u(t,x)\cdot v+
            \tfrac12\theta(t,x)(|v|^2-5),
\]
where $(u,\theta)$ solves an incompressible
Navier--Stokes--Fourier system and satisfies its energy inequality
\citep[Theorem 2.4]{microGolseSaintRaymond2009}.
The initial velocity is determined by the weak limit of the Leray
projection of the scaled momentum; a corresponding thermal moment
determines the initial temperature fluctuation.  For appropriate
well-prepared initial entropy, the velocity is a Leray solution
with the corresponding initial energy bound.

This is a theorem about a specified collision class, spatial domain,
scaling, solution concept, and convergence topology.  Its proof uses
entropy bounds, coercivity of the linearized collision operator,
velocity averaging, vanishing conservation defects, and control of
nonlinear fluxes.  It yields weak fluid solutions; it does not by
itself upgrade them to every stronger regularity or uniqueness
statement.  Statements about the separate 2026 proof release are
analyzed in Chapter~\ref{ch:navier-stokes-release}.  That chapter's
precise theorem must be used with its own forcing and solution
hypotheses, rather than inferred from this historical limit theorem.

\section{DSMC as a stochastic approximation, with auditable assumptions}

The requested sixth Volkov deck explains direct simulation Monte Carlo
as a particle approximation to kinetic dynamics, with statistical
weights, free-flight and collision substeps, cell-based collision
selection, inflow and wall kernels, and sampling of moments
\citep{microVolkov06}.  These simulated particles represent weighted
samples.  Their paths are not claimed to be the paths of every physical
molecule in the gas.

Three error categories must be distinguished: the physical model
(collision law and wall interaction), numerical discretization
(time step, cells, particle count, cutoffs), and statistical estimation.
Increasing a GPU sample count primarily addresses the latter two
categories.  It cannot validate a physical collision law by itself.

\begin{proposition}[Flux sampling at an equilibrium wall]
\label{prop:micro-flux}
Let a stationary wall emit an equilibrium gas with thermal variance
$\theta>0$ and inward unit normal $n$.  The normalized velocity
density of molecules crossing the wall into the domain is
\[
 p_{\mathrm{flux}}(v)
 =\frac{(v\cdot n)_+}{2\pi\theta^2}
       e^{-|v|^2/(2\theta)}.
\]
Its positive normal component $s$ has density
$(s/\theta)e^{-s^2/(2\theta)}$, and its two tangential
components are independent centered Gaussians of variance $\theta$.
\end{proposition}
\begin{proof}
A bulk density $M_\theta(v)$ contributes a crossing flux
$(v\cdot n)_+M_\theta(v)$.  The integral of this flux is
$\sqrt{\theta/(2\pi)}$.  Dividing by that value gives the
displayed density.  In orthonormal normal/tangential coordinates
the exponential and its prefactor factor into the three stated
one-dimensional densities.
\end{proof}

Sampling the normal velocity from an unweighted half-Gaussian would
therefore give the wrong boundary ensemble.  A finite reservoir
implementation also requires an estimate of the missing high-speed
tail; the Maxwellian has no finite maximum speed.

\begin{proposition}[The exact frozen-rate majorant identity]
\label{prop:micro-majorant}
Consider $K\ge1$ possible events with fixed rates $r_j\ge0$, and
choose $R>0$ with $R\ge\max_jr_j$.  Generate a Poisson proposal process of rate
$KR$, choose a label uniformly from $\{1,\ldots,K\}$ at each
proposal, and accept label $j$ with probability $r_j/R$.
The accepted events of label $j$ form a Poisson process of rate
$r_j$; distinct accepted labels have independent counts.
\end{proposition}
\begin{proof}
The probability that a proposal is accepted with label $j$ is
$r_j/(KR)$.  For an interval of length $t$, condition on its
Poisson number of proposals and take the joint probability
generating function of the marked counts.  Summing the Poisson
series gives
\[
 \exp\!\left(t\sum_{j=1}^K r_j(z_j-1)\right)
 =\prod_{j=1}^K\exp(tr_j(z_j-1)).
\]
This is the product of the generating functions of independent
Poisson variables with the claimed means.  The same argument on
disjoint time intervals proves the process statement.
\end{proof}

This proves the sampling principle behind majorant collision schemes.
It does not alone prove convergence of a full DSMC method whose rates
change after collisions and whose transport is split in time.  In that
case the bound must remain valid for the current state and the
discretization error must be controlled.  If an alleged acceptance
probability exceeds one, silently clipping it changes the rate.
A temperature-based constant in illustrative code is not a proof of
a global bound on an unbounded velocity distribution.

\begin{proposition}[Sampling variance with correlations]
\label{prop:micro-sampling}
If $Y_1,\ldots,Y_n$ have finite second moments, then
\[
 \operatorname{Var}\!\left(\frac1n\sum_iY_i\right)
 =\frac1{n^2}\sum_{i,j}\operatorname{Cov}(Y_i,Y_j).
\]
For identically distributed independent samples this equals
$\sigma^2/n$.  For a stationary sequence with variance
$\sigma^2>0$ and lag correlations $r_k$, it equals
\[
 \frac{\sigma^2}{n}
 \left(1+2\sum_{k=1}^{n-1}(1-k/n)r_k\right).
\]
\end{proposition}
\begin{proof}
Expand the square of the centered sum and take expectations.
For a stationary sequence, group off-diagonal pairs by their
lag; there are $n-k$ pairs at lag $k$ in each direction.
\end{proof}

Thinning a time series can reduce some correlations but does not prove
independence.  A central-limit error bar requires its own assumptions;
a fixed seed provides reproducibility, not randomness certification.
For independent samples, Chebyshev's inequality gives the explicit
bound
$\Pr(|\overline Y-\mathbb EY|\ge a)\le\sigma^2/(na^2)$.
The statistics and Monte Carlo decks are useful introductions
\citep{microVolkov04,microVolkov05}; implementations should additionally
handle finite precision, zero arguments of logarithms, array bounds,
correct random-number-generator state, and reproducible restart state.
The period of a generator is constrained by its internal state space,
not merely by the number of possible values of a single output word.

\section{Reduced models and learned closures}

An orthogonal reduced basis minimizes a chosen projection loss on its
training snapshots.  That statement does not imply that the projected
trajectory is stable, that the discarded modes exert no force, or that
the basis will predict new parameter regimes.  The currently available
replacement Farhat course decks explicitly distinguish projection
error from model-dynamics error \citep{microFarhatCourse}.  Their
complete version-specific review is recorded separately in the source
audit.

For an orthonormal column matrix $V$, write a full state
$z=Va+w$ with $V^Tw=0$.  If $\dot z=F(z)$, the exact retained
coordinate satisfies $\dot a=V^TF(Va+w)$.  Replacing the right
side by $V^TF(Va)$ is a closure.  The exact memory identity in
Chapter~\ref{ch:quantum-fluid} explains why eliminated coordinates
can leave both memory and an initial-condition contribution.

\begin{proposition}[What weighted POD actually minimizes]
\label{prop:micro-pod}
Let $X=[x_1\ \cdots\ x_s]\in\R^{n\times s}$, let
$W=\operatorname{diag}(w_1,\ldots,w_s)$ with $w_j>0$, and
let $\lambda_1\ge\cdots\ge\lambda_n\ge0$ be the eigenvalues
of $C=XWX^T$.  Fix $1\le r\le n$.  Among all
$V\in\R^{n\times r}$ with $V^TV=I$,
\[
 \min_V\sum_jw_j\|x_j-VV^Tx_j\|^2
 =\sum_{i>r}\lambda_i.
\]
A minimizing range is spanned by eigenvectors for the largest
$r$ eigenvalues.
\end{proposition}
\begin{proof}
The objective is $\tr C-\tr(V^TCV)$.  In an orthonormal
eigenbasis $e_i$ of $C$, put $a_i=\|V^Te_i\|^2$.
Then $0\le a_i\le1$, $\sum_i a_i=r$, and
$\tr(V^TCV)=\sum_i\lambda_i a_i$.  Transferring any weight
from a smaller eigenvalue to an unfilled larger one cannot decrease
this sum.  Its maximum is therefore $\sum_{i\le r}\lambda_i$,
and the indicated eigenspace attains it.
\end{proof}

The proposition concerns the chosen snapshots and weights.  It contains
no assertion about an unsampled parameter or a differential equation.
Indeed the stable matrix
\[
 A=\begin{pmatrix}-1&10\\0&-1\end{pmatrix},\qquad
 V=2^{-1/2}\begin{pmatrix}1\\1\end{pmatrix}
\]
has both eigenvalues equal to $-1$, while its orthogonal Galerkin
reduction is $V^TAV=4$.  The reduced scalar solution grows
exponentially.  Nonsymmetric stable dynamics are therefore not
automatically preserved by orthogonal projection.

For Petrov--Galerkin projection with full-column-rank trial and test
matrices $V,W$, the matrix $W^TV$ must be invertible to use
$\Pi=V(W^TV)^{-1}W^T$.  Direct multiplication gives
$\Pi^2=\Pi$, $\operatorname{range}\Pi=\operatorname{range}V$,
and $\ker\Pi=\ker W^T=(\operatorname{range}W)^\perp$.
The direction of this oblique projection is the latter kernel.

\subsection{Hyperreduction adds another controlled approximation}

State reduction need not make a nonlinear residual cheap to evaluate.
Let $U\in\R^{n\times r}$ have orthonormal columns and let
$P\in\R^{n\times r}$ select $r$ coordinates.  If
$P^TU$ is invertible, the discrete empirical interpolation map is
\[
 D=U(P^TU)^{-1}P^T.
\]

\begin{proposition}[Interpolation exactness and amplification]
\label{prop:micro-deim}
For every $y\in\R^n$, $DU=U$, $P^TDy=P^Ty$, and
\[
 \|y-Dy\|
 \le\left(1+\|(P^TU)^{-1}\|\right)
       \|(I-UU^T)y\|.
\]
\end{proposition}
\begin{proof}
The first two identities follow by cancellation of $P^TU$.
Because $D$ is the identity on the range of $U$,
$(I-D)y=(I-D)(I-UU^T)y$.
Use $\|I-D\|\le1+\|D\|$ and the unit operator norms
of $U$ and $P^T$ to obtain the estimate.
\end{proof}

The bound is deliberately elementary; sharper oblique-projection
bounds are possible.  Even this version shows why a small discarded
snapshot norm is insufficient if the sampled basis is ill-conditioned.
For oversampled least squares, full column rank of $P^TU$ is the
corresponding uniqueness and exact-reproduction requirement.
An online cost independent of $n$ additionally requires that sampled
nonlinear entries and the restricted reconstruction can be evaluated
without assembling all $n$ entries.  A dense nonlocal nonlinear map
can violate that requirement.

Other reduction methods have different hypotheses.  Classical balanced
truncation concerns stable finite-dimensional linear time-invariant
systems, with reachability and observability conditions, and its usual
input--output error estimate assumes zero initial state.  An empirical
balanced-POD approximation does not inherit every exact balancing
guarantee.  Moment matching requires the relevant full and reduced
resolvents to exist at its interpolation nodes.  Grassmann-chart
interpolation requires the chart to be valid, fixed subspace rank,
and appropriate alignment before interpolating reduced operators.
The detailed Farhat review documents these distinctions
\citep{microFarhatCourse}.

Likewise, a weighted finite-element energy
$\widetilde U(q)=\sum_e\xi_e U_e(q)$ with fixed weights has
gradient $\sum_e\xi_e\nabla U_e(q)$.  If a neural model makes
$\xi_e=\xi_e(q)$, the gradient instead includes
$\sum_e U_e(q)\nabla\xi_e(q)$.  Omitting those terms loses the
claimed potential-force relation.  A preservation theorem for a
fixed-weight mechanical hyperreduction must therefore be re-proved
when the weights become state dependent or the equation changes to
a different fluid formulation.

\subsection{A provable constraint for a neural constitutive model}

A classical or quantum neural model can be used to predict constitutive
coefficients from invariant local features.  For example, a model can
output nonnegative scalars $\mu,\zeta,\kappa$ in
\eqref{eq:micro-constitutive}.  This architecture has a limited but
useful exact guarantee.

\begin{proposition}[Nonnegative constitutive entropy production]
\label{prop:micro-learned-entropy}
For $T>0$, let the viscous stress be
$\tau=2\mu S_0+\zeta(\nabla\cdot u)I$ and the heat flux
$q=-\kappa\nabla T$, with $\mu,\zeta,\kappa\ge0$.
Then
\[
 \frac{\tau:\nabla u}{T}
       -\frac{q\cdot\nabla T}{T^2}
 =\frac{2\mu|S_0|^2}{T}
  +\frac{\zeta(\nabla\cdot u)^2}{T}
  +\frac{\kappa|\nabla T|^2}{T^2}\ge0.
\]
\end{proposition}
\begin{proof}
The symmetric stress is orthogonal in the Frobenius inner product to
the antisymmetric part of $\nabla u$.  The traceless symmetric
part $S_0$ is orthogonal to $I$.  Decompose the symmetric part of
$\nabla u$ into $S_0+(\nabla\cdot u)I/3$ and substitute the
heat-flux formula.
\end{proof}

If the scalar inputs are rotational invariants, these constitutive
tensors transform correctly under rotations.  Dependence on gradients
and thermodynamic invariants also permits Galilean covariance.
Neither property proves accurate transport coefficients.  Moreover,
nonnegative entropy production does not by itself make
$P=pI-\tau$ positive semidefinite far from equilibrium.  A kinetic
moment model that requires realizable covariances can instead
parameterize $P=LL^T$ or another positive-definite form and impose
the remaining energy and flux constraints explicitly.  These are
different requirements and must be checked together.

Tensor decomposition makes these constraints transparent.  In three
dimensions a rank-two tensor splits into scalar trace, antisymmetric,
and symmetric traceless pieces, of dimensions $1,3,5$.  The
antisymmetric piece is an exterior-square representation; the
symmetric pieces arise from symmetric tensor powers and their
trace decomposition.  Schur functors and invariant contractions
therefore help build features and outputs with known transformation
laws.  Determinants test nonsingularity or volume change in specified
finite-dimensional maps.  They do not determine a turbulent closure
or imply the analytic estimates required for a continuum theorem.

For a quantum version of ToricGT, the surface or dessin specifies
connectivity and admissible local operations, while circuit parameters
specify a family of quantum channels and observables.  A conservation
law can be enforced by restricting channels or observables to a
specified invariant sector.  A measured expectation then estimates a
well-defined feature or coefficient.  State preparation, finite-shot
variance, circuit approximation, and the relation between that feature
and the physical kinetic state remain separate parts of the model.
The quantum-fluid chapter discusses these algorithmic costs and
encodings.  Learning on a surface does not identify an arbitrary
three-dimensional fluid with a two-dimensional toric code.

\section{What passes from complexes to a continuum theorem}

On a finite cell complex, the boundary-of-boundary identity supplies
exact incidence relations.  Its adjoint cochain identities, Hodge
decompositions, and positive semidefinite Laplacians can enforce
discrete conservation and isolate harmonic components.  A gentle
surface algebra or a dessin provides additional combinatorial
organization of local operators.  A variety of complexes organizes
matrices with prescribed compositions and rank conditions.  These
constructions can certify the algebraic structure of a discretization
or constrain a neural architecture.

To pass from such a model to a continuum claim requires a separate
chain of estimates: consistency of the geometric approximation,
stability uniform in mesh size, compactness in a topology suitable
for nonlinear products, identification of the limiting equation,
and preservation of the required initial and boundary data.
Spectral gaps on individual finite graphs need not have a uniform
positive lower bound as the graphs are refined.  Rank constraints
on finitely many matrices need not control high-frequency
concentration.  A learned closure with positive entropy production
need not supply the required continuum regularity.

The constructive role of the project's algebra is to make some
requirements exact and testable, then to identify which analytic
requirements are still independent.  Chapter~\ref{ch:fluid-structures}
proves the finite structural statements, and
Chapter~\ref{ch:navier-stokes-release} identifies the actual theorem
and proof obligations of the recent release.  The microscopic
derivations in this chapter provide the physical interpretation and
the correct hierarchy of approximations without changing those
theorems' hypotheses.
